%================================================================
% CHAPTER 1: INTRODUCTION
%================================================================
\chapter{INTRODUCTION}

\section{Introduction}

This chapter contains the introduction of the project that is, problem statement, problem domain of the project, project features or modules, hardware requirements and software requirements of the project which are used for development of the proposed system and how to access or use this system. This chapter also provides a discussion of project overview with comprehensive view. Problem domain covers the problems exist in the current system and problem statement describes the problem which is solved by the project. Objective describes the explanation about what is the aim of designing and developing this application. Project features give the brief description about the quality of the project. Project scope describes the scope or area of the project.

\section{Project Overview}

This project is a web-based project in order to increase efficiency, maintainability and managing information of universities. This site is developed to minimize the gap through visiting one website provided by Academia Portal. We intend to make a platform where a user can get all the information by selecting the degree programs and location at where they want to study. This proposed project extracts data from the websites of different universities and accumulates on one website. In this project we have extracted information about different universities of twin cities which includes program offered, fee structure, university ranking, eligibility criteria, faculty, admissions date and entry tests sample papers. Users do not need to visit all the websites individually; they can find all the university information at one platform. Another purpose of this project is to compare the different universities on different criteria including fee structure and HEC ranking.

In this proposed project there is also a choice for searching through degree program or course. Academia Portal also displays a list of universities which offer the degree program. The system assures the availability of the right information at the right time. This proposed system makes sure the availability of a backup system, which means if any university website gets down and a user cannot see the data of that university website. Academia Portal gives the data from our database, that is made available by collecting data and crawling data from different websites.

\section{Problem Domain}

Many internet websites are today working not as automated; like each university has its own website in which they show all the information about their different criteria. Users have to individually checkout all websites one by one manually and then get required information. In Pakistan there is no effective and efficient platform or website which works in the way that can compare and do other many functionalities on a single website. Another problem is that there is no solution to the failure of any university website. After that, a user is not able to get information in the case of failure. So it is going to be worst for the user because maybe that university is good and preferable but it gets down due to any failure.

\section{Problem Statement}

In these days many websites share their information with other universities with their own websites. A single website takes much time; users need to visit websites individually one by one which is very time consuming, and no one can understand easily how to get information what he or she requires. There is not any single website that compares university data with each other. This proposed system helps to provide all information on one platform about different universities to students or users in an efficient way and compare data like fee, ranking, faculty information. It is then very easy for users to get information that they want, and also the proposed system works 24/7 on the internet. This system can be used by fresh students and other users who take admission in university. After accessing the system they can be authorized or allowed to view different information of different universities. All the information about websites is stored in the database of the system after extraction and filtration. As it is mentioned in the problem domain, within the case of failure the user cannot get the data that they want, so to overcome this difficulty the best way to solve this problem is the backup system.

\section{Aims and Objectives}

The main objective of developing and implementing this project is to collaborate with students. It provides a fast and efficient way to analyze university information on one website. This system reports to users at any time. It can update all the updated information and data which are newly added by any university website. Recently most of the students visit websites individually which takes a lot of time and many websites remain unchecked. Thus, the focus of this project is to analyze websites in a fast and automated way and give complete and perfect results within no time. We will also have to provide the academic and educational information of different other universities of different countries. It has become furnished further on the user demands.

\section{Project Features}

Basically the proposed system is used to resolve the problem and has to give a more effective way. This system is just developed for the twin cities of Pakistan. This system is very comfortable in manner of usage for all kinds of users. Features of the proposed system are described below.

\subsection{Search Query}

In this module the user should be able to search different universities which exist in the twin cities of Pakistan. The user can also search different universities on the basis of program by selecting a specific or entire program; then the system displays the list of those universities which are offering the entire program.

\subsection{General Information}

In this feature, users search general information of different universities. The features covered in this module are given below.

\subsubsection{Fee Structure Information}

In this module the system provides all the details of the fee structure. Users should be able to view the fee structure of all the courses of all the universities which are included in this project. In this module the user is able to search or view the fee structure of different universities by selecting university name, university department and the course category.

\subsubsection{Program Offered Information}

In this module the system contains all the information about programs offered in different universities. Users should be able to view or search about different programs offered in different universities. Users should search programs by selecting department, category and course; then the response system gives the output in table form showing that program is offered by this university.

\subsubsection{Faculty Information}

This feature contains all universities' faculty information and provides it to users. Users can also be able to obtain information about faculties of different departments of different universities.

\subsubsection{Department Information}

This feature contains all information about departments and provides it to users. Through this, users can get all information about departments of different universities.

\subsubsection{Eligibility Criteria}

This proposed project also extracts information about eligibility criteria of all the universities. Through this, users can also be able to view all the universities' eligibility criteria by selecting university name and course without visiting all the university websites.

\subsubsection{Ranking}

In this module the proposed project also extracts information about the general ranking of all the universities which are available on the Higher Education Commission (HEC) website.

\subsection{Comparison}

This feature defines or describes different criteria on which we can compare different universities. In this feature we can compare different universities on the basis of fee structure and on the basis of HEC-based ranking. Users should be able to compare fee structure of a single course or program and also should be able to compare all the courses of all the universities on one page just on one click.

\section{Resource Requirements}

There are two types of resources which are required for developing the system.

\subsection{Software Requirements}

The software resources required for our system are Windows 7, Windows 8, 8.1 and Windows 10, Windows XP, and PHP.

For developing purpose, we need:

\begin{itemize}
    \item PHP, Microsoft SQL Server, Windows XP, Windows 8 and above.
\end{itemize}

\subsection{Hardware Requirements}

The hardware required for the system is a dual core or latest processor. The device on which this system runs should have a minimum 1 gigahertz (GHz) or faster processor and at least 2 gigabytes (GB) of RAM. The system must have up to 40 GB hard disk space and also have an internet connection.

\section{Tools and Technologies}

For this project we have used the updated and efficient versions of tools for the development of this proposed system. The system uses PHP as it provides fast and efficient development of the system. For maintaining the database, the system uses SQL Management Studio 2012. Details are as follows.

\subsection{PHP}

The Personal Home Page Tools, also called PHP, is the strongest development tool that began in 1995. PHP is a server-side scripting language designed primarily for web development but is also utilized as a broadly useful programming language. PHP can build simple, dynamic and static web applications. It is used to quicken bug reporting and to enhance the code.

\subsection{CSS}

CSS stands for Cascading Style Sheets. Cascading Style Sheets is a style sheet language that is used to make the better look of a markup language written document, meaning HTML. CSS is designed for the purpose of separating the content and presentation of the web page, which includes colors, fonts and layout. This separation helps in improving the accessibility of the content, more flexibility and in controlling specification of the presentation characteristics. There are two styles of CSS used in designing web pages: internal CSS and external CSS. Internal CSS is applied within an individual HTML tag to use change only for that line, paragraph, form, etc. External CSS is that in which developers make an external file of CSS in a folder of CSS. It enables developers to only give a reference of that CSS file and it will change the layout of the complete HTML file. External CSS reduces the complexity and repetition of patterns in the content structure.

\subsection{JavaScript}

JavaScript is abbreviated often as JS. JavaScript is an interpreted high-level programming language. JavaScript is a language which is also considered multi-paradigm, weakly typed, dynamic and prototype-based. Like HTML and CSS, JavaScript is one of the three core technologies of the World Wide Web. JavaScript is an essential element in developing web-based programs that can assist in creating interactive web pages. JavaScript is used in a wide range of websites and a dedicated JavaScript engine is available on all the major web browsers to execute it. As JavaScript is a multi-paradigm language, it supports functional, event-driven and imperative programming patterns. JavaScript has an API (Application Programming Interface) for operating with arrays, text, dates, regular expressions and also some basic influence of DOM (Document Object Model).

\subsection{HTML}

HTML stands for Hypertext Markup Language. HTML is the language used by developers for creating distinct websites, web pages and applications on the internet. These websites, web pages and applications are designed with the help of other technologies which are used to create stunning and attractive content. Those technologies are Cascading Style Sheets (CSS) and JavaScript. Web browsers receive files of HTML from the server of internet or from any local storage. HTML explains the structure of the web page semantically and initially included indications for the document's appearance. Elements of HTML are used as building blocks for HTML and web pages. HTML consists of images, paragraph writing, making or growing interactive forms and creating tables.

\subsection{Bootstrap}

Bootstrap is an open source free front-end framework. It is used for faster and simpler web development. It includes HTML and CSS-based design templates for user interface components like forms, tables, buttons, dropdowns, navigations, and modals, indicators, tabs, carousel and several other optional JavaScript extensions. It provides the power to create a responsive layout with less effort. The main advantage of using Bootstrap is that it comes with a free set of tools for building responsive and flexible web layouts as well as common interface components.

\subsection{Microsoft SQL Server}

Microsoft SQL Server Management Studio is comprehensive, integrated data management and analysis software that enables organizations to reliably manage mission-critical information and confidently run today's increasingly complex business applications. SQL Server allows companies to gain greater insight from their business information and achieve faster results for a competitive advantage.

\section{Project Scope}

The scope of the proposed project is that it only extracts the websites data with domain ``\texttt{.pk}'' and ``\texttt{.com}''. The proposed system extracts the contents of webpages using Crawler, scraper and DOM parsing library to extract only ``textual information'' from different web pages or sites and stores in the database after filtration. This system cannot extract websites with extension like ``\texttt{.edu}'', ``\texttt{.org}'' etc. The system allows specific users to get specific information about different university websites on one page.

\section{Report Layout}

The outline of the project report is described below.

The first chapter describes the introduction of the project. It includes problem domain, problem statement, development methodology, aims and objectives, project scope, tools and technologies and resource requirements. The second chapter includes a discussion on literature review. It includes background, existing work, methodologies used for content mining and web mining used in data extraction, data filtration methods and project limitations. The third chapter explains requirement specification. Functional and non-functional requirements and modules of the system are explained in detail. The next chapter consists of proposed system and system modeling: system analysis, evaluation of existing work, proposed system, use case and use case diagram, sequence system with diagram, data flow diagram, and architecture diagram. The fifth chapter describes the result, testing, validation and analysis of the system and different test modules and different testing techniques. The sixth chapter consists of conclusion, limitations and future work.
